\subsection{Метод Гаусса решения линейных систем. Теорема Кронекера-Капелли}

\begin{definition}
Элементарными преобразованиями системы линейных уравнений называются следующие действия:
\begin{enumerate}
    \item Перестановка двух уравнений системы.
    \item Умножение уравнения на число $\lambda \neq 0$.
    \item Прибавление к одному из уравнений другого уравнения, умноженного на произвольное число.
\end{enumerate}
Данные преобразования соответствуют элементарным преобразованиям строк расширенной матрицы системы.
\end{definition}

\begin{theorem}
После элементарных преобразований системы уравнений получится система, эквивалентная исходной (имеющая тот же набор решений).
\end{theorem}

\begin{tcolorbox}[title=Доказательство, breakable]
\textbf{1. Перестановка уравнений.} Очевидно, что от перемены порядка уравнений набор решений не меняется.

\textbf{2. Умножение на $\lambda \neq 0$.} Если набор чисел является решением уравнения, то при подстановке в него получается верное числовое тождество. При умножении тождества на ненулевое число оно остается верным тождеством. В силу обратимости (деление на $\lambda$) решение нового уравнения будет решением старого.

\textbf{3. Прибавление уравнения к уравнению.} Пусть к первому уравнению $\sum a_{1j}x_j = b_1$ прибавили второе $\sum a_{2j}x_j = b_2$, умноженное на $\alpha$. Если $(\gamma_1, \dots, \gamma_n)$ — решение исходной системы, то выполняются тождества:
\[ \sum a_{1j}\gamma_j = b_1, \quad \sum a_{2j}\gamma_j = b_2 \]
Складывая их (второе с коэффициентом $\alpha$), получаем:
\[ \sum (a_{1j} + \alpha a_{2j})\gamma_j = b_1 + \alpha b_2 \]
что является первым уравнением новой системы. Таким образом, всякое решение исходной системы является решением новой. В силу обратимости элементарных преобразований (вычитание второй строки, умноженной на $\alpha$, из первой), всякое решение новой системы является решением старой.
\end{tcolorbox}

\begin{definition}
Матрица называется \textbf{ступенчатой}, если выполнены два условия:
\begin{enumerate}
    \item Все нулевые строки (если они есть) стоят последними.
    \item Для каждой ненулевой строки её первый ненулевой элемент (лидер) расположен строго правее лидера предыдущей строки.
\end{enumerate}
\end{definition}

\begin{theorem}
Всякую ненулевую матрицу путём элементарных преобразований только над строками можно привести к ступенчатому виду.
\end{theorem}

\textbf{Алгоритм метода Гаусса:}
\begin{enumerate}
    \item \textbf{Прямой ход:} С помощью элементарных преобразований расширенная матрица системы $(A|B)$ приводится к ступенчатому виду. Если на этом этапе возникает строка вида $(0, 0, \dots, 0 | b^*)$, где $b^* \neq 0$, то система \textbf{несовместна} (противоречие $0 = b^*$).
    \item \textbf{Анализ ступенчатого вида:} Неизвестные, соответствующие столбцам с лидерами строк, называются \textbf{главными}, остальные — \textbf{свободными}. Количество главных переменных равно рангу матрицы.
    \item \textbf{Обратный ход:} Свободным переменным придаются произвольные значения (параметры). Главные переменные выражаются через свободные (начиная с последнего уравнения к первому).
\end{enumerate}

\begin{theorem}[Кронекера--Капелли]
Система линейных алгебраических уравнений совместна тогда и только тогда, когда ранг основной матрицы системы $A$ равен рангу расширенной матрицы системы $\tilde{A}$:
\[ \mathrm{rk}(A) = \mathrm{rk}(\tilde{A}) \]
\end{theorem}

\begin{tcolorbox}[title=Доказательство, breakable]
Приведем систему к ступенчатому виду методом Гаусса. Обозначим ступенчатые формы основной и расширенной матриц через $A_{ступ}$ и $\tilde{A}_{ступ}$. При элементарных преобразованиях ранг матрицы не меняется, поэтому $\mathrm{rk}(A) = \mathrm{rk}(A_{ступ})$ и $\mathrm{rk}(\tilde{A}) = \mathrm{rk}(\tilde{A}_{ступ})$.

$\implies$ Пусть система совместна. В методе Гаусса это означает отсутствие противоречивых уравнений вида $0 = b^*$ при $b^* \neq 0$. Следовательно, в ступенчатом виде расширенной матрицы нет строк, где лидер стоит в столбце свободных членов. Это означает, что количество ненулевых строк в $A_{ступ}$ и $\tilde{A}_{ступ}$ одинаково. Так как ранг ступенчатой матрицы равен числу её ненулевых строк, получаем $\mathrm{rk}(A) = \mathrm{rk}(\tilde{A})$.

$\impliedby$ Пусть $\mathrm{rk}(A) = \mathrm{rk}(\tilde{A}) = R$. Это означает, что количество ненулевых строк в ступенчатых формах основной и расширенной матриц совпадает. Следовательно, в расширенной матрице не может появиться строка с лидером в последнем столбце (столбце свободных членов), так как тогда ранг расширенной матрицы был бы на единицу больше ранга основной. Отсутствие таких строк гарантирует по методу Гаусса возможность нахождения решения, то есть совместность системы.
\end{tcolorbox}

\textbf{Следствия из теоремы Кронекера--Капелли:}
\begin{enumerate}
    \item Число главных переменных в совместной системе равно её рангу $R$.
    \item Совместная система является \textbf{определенной} (единственное решение) тогда и только тогда, когда ранг системы равен числу неизвестных ($R = n$).
    \item Совместная система является \textbf{неопределенной} (бесконечно много решений), если ранг меньше числа неизвестных ($R < n$).
\end{enumerate}

\begin{tcolorbox}[title=Источники, colback=gray!10]
\begin{itemize}
  \item Лекция №\,2, тема «Метод Гаусса решения линейных систем», временной отрезок 67:00--90:00
  \item Лекция №\,3, тема «Теорема Кронекера-Капелли», временной отрезок 00:00--05:00
  \item PDF «Линейная алгебра\_compressed (1).pdf», разделы 3 и 4 по СЛАУ, стр. 10--13
  \item PDF «Вопросы», вопрос №\,5
\end{itemize}
\end{tcolorbox}
